\documentclass[10pt,a4paper]{amsart}
\usepackage[margin=19mm]{geometry}
\usepackage{amsmath,amssymb,mathtools}
\usepackage[scr=rsfs,cal=euler]{mathalfa}
\usepackage{xfrac}
\usepackage[unicode,hidelinks,bookmarks=false]{hyperref}
\newcommand{\ie}{i.e.\ }
\newcommand{\cf}{cf.\ }
\newcommand{\resp}{respectively\ }
\newcommand{\Subsection}{\subsection}
\providecommand{\nobreakdash}{\nobreak\hbox{-}}
\setlength{\emergencystretch}{2em}
\multlinegap0pt
\usepackage{mathtools}
\usepackage{amssymb}
\usepackage{enumerate}
\newtheorem{theorem}{Theorem}
\newcommand{\depth}{\operatorname{depth}}
\newcommand{\image}{\operatorname{image}}
\newcommand{\ov}{\overline}
\newcommand{\rt}{\rightarrow}
\title{Annihilators of local cohomology modules over modular invariant rings and Dickson polynomials}
\author{Tony J. Puthenpurakal}
\date{Source: arXiv:2603.16491v2 (CC BY 4.0)}
\begin{document}
\maketitle

\noindent\emph{Source and licence.} Annihilators of local cohomology modules over modular invariant rings and Dickson polynomials. Version arXiv:2603.16491v2 (CC BY 4.0), distributed by arXiv under the Creative Commons Attribution 4.0 International licence, http://creativecommons.org/licenses/by/4.0/, as recorded in the arXiv licence metadata for this exact version and shown on the abstract page https://arxiv.org/abs/2603.16491v2. Full licence notice: \texttt{licenses/arxiv-2603.16491-CC-BY-4.0.txt}.\par\smallskip

\noindent\textit{Editorial scope.} Counted unit: the article's theorem dual-dim (source0.tex lines 645--647). The article's complete original proof of it is delivered with it (source0.tex lines 651--668, one \texttt{proof} environment of 18 lines, which the article places away from the statement). No mathematics is written, repaired or re-derived: the statement and the proof are the article's own lines, in the article's own notation, and no retained line is substituted or re-flowed.  No further source-local context is needed: every cross-reference of every retained line resolves inside the two delivered spans.  Nothing was omitted from any retained span, so the package carries no omission record.  No retained line contains a citation, so the wrapper carries no bibliography and no bibliography entry.  No retained line contains a figure environment, an image-inclusion command or drawing code, so no image is delivered and the package declares no asset dependency.  Editorial formatting only, in addition: an AMS \texttt{amsart} preamble that restates the article's own declarations and macros used by the retained lines, namely the environments \texttt{theorem} on separate counters, together with the standard packages those lines need.  The retained environments print in this document as Theorem 1 in order of appearance, whereas the article prints the same items with the numbering of its own full document.  The complete native source of the article as distributed in the arXiv e-print for this exact version is bundled unchanged as \texttt{sources/196555/source0.tex}.
\begin{theorem}
\label{dual-dim} Let $K$ be a field. Let $T= \bigoplus_{n \geq 0}T_n$ be a finitely generated graded $K = T_0$-algebra. Let $M$ be a finitely generated graded $T$-module of dimension $r$. Let $U_i = H^i_{S_+}(M)$  for $i \geq 0$. Note $U_i^{\vee}$ is a finitely generated graded $T$-module. Then $\dim U_i^\vee \leq i$ for $0 \leq i \leq r$.
\end{theorem}

\begin{proof}[Proof of Theorem \ref{dual-dim}]
Set $H^i(-)  = H^i_{T_+}(-)$.
We prove by induction on $i$ that if $M$ is a finitely generated graded $T$-module of dimension $s$ then $H^i(M)^\vee $ has dimension $\leq i$ for $i \geq 0$.
This result is certainly true when $i = 0$.
Assume $i \geq 1$ and that $\dim H^j(N) \leq j$ for all finitely generated graded $T$-modules with $j < i$.

If $E$ is a finitely generated, graded $T$-module then set $\ov{E} = E/H^0(E)$. If $\dim E > 0$ then $\depth \ov{E} > 0$ and $H^r(\ov{E}) = H^r(E)$ for $r > 0$. Let $x \in T_+$ be homogeneous such that $x$ is $\ov{E}$-regular. Then $\ker (E(-|x|) \xrightarrow{x} E)$ has finite length.

If $\dim M = 0$ then $H^r(M) = 0$ for $r > 0$. So we have nothing to show. Assume $\dim M > 0$.
Also if $\dim H^i(M)^\vee \leq 0$ then again  we have nothing to show. So assume $\dim H^i(M)^\vee > 0$. Let $x \in T_+$ be homogeneous which is $\ov{M} \oplus \ov{H^i(M)^\vee}$-regular.
We note that $H^j(M) \cong H^j(\ov{M})$ for all $j \geq 1$.
The exact sequence $0 \rt \ov{M}(-|x|) \xrightarrow{x} \ov{M} \rt N \rt 0$ induces a long
exact sequence
$$H^{i-1}(N) \xrightarrow{\delta} H^i(M)(-|x|) \xrightarrow{x} H^i(M).$$
Set $W = \image \delta$. Taking duals we get an exact sequence
$$ 0 \rt L \rt H^i(M)^\vee \xrightarrow{x} H^i(M)^\vee(|x|) \rt W^\vee \rt 0$$
and that $W^\vee $ is a submodule of $H^{i-1}(N)^\vee$ and so by induction hypothesis  has dimension $\leq i -1$. As $L$ has finite length it follows that $\dim H^i(M)^\vee \leq i$. The result follows by induction.
\end{proof}

\end{document}
